\appendix

\section{Data Preparation Details}
\label{app:data-prep}

\paragraph{Source and unified tables.}
The extraction consumed the OMOP tables \texttt{person},
\texttt{visit\_occurrence}, \texttt{visit\_detail}, \texttt{death},
\texttt{condition\_occurrence}, \texttt{observation},
\texttt{procedure\_occurrence}, \texttt{device\_exposure},
\texttt{drug\_exposure}, and \texttt{measurement}, and reorganized them into
five unified tables: \emph{visits}, \emph{visit types}, \emph{clinical
events} (conditions, observations, procedures, and device exposures under a
common schema), \emph{drug exposures}, and \emph{measurements}. The only
enrichments applied were the attachment of Anatomical Therapeutic Chemical
(ATC) codes to drug exposures, where available, the harmonization of column
types, and the closing of intervals whose end timestamp was missing or not
later than their start one second after their start, so that point
observations stored under an interval schema carry an explicit rather than a
degenerate duration. Dimension tables enumerating the distinct clinical,
measurement, and drug concepts, together with their units and administration
routes, were derived from the unified tables and used to support concept
curation.

\paragraph{Concept curation.}
Candidate codes for each target concept were surfaced by descriptive
analytics over the dimension tables, ranked by frequency of occurrence and
by apparent relevance to the concept, and then accepted or rejected by the
clinical domain expert. The resulting indices define both the concepts
included in the abstraction input and the departmental and diagnostic
exclusions applied in the cohort definition.

\paragraph{Derived concepts.}
Body mass index was computed from paired height and weight measurements
where not recorded directly. Ketoacidosis events were derived from
combinations of pH, bicarbonate, and intravenous insulin administration, as
described in Appendix~\ref{app:definitions}.

\paragraph{Modified Charlson comorbidity index.}
The score is a weighted sum of ten chronic condition flags ascertained from
the first 48 hours of admission: malignancy, chronic kidney disease, and
chronic liver disease weighted 2; and cardiac disease, chronic respiratory
disease, diabetes, dementia, alcohol misuse, anemia, and recurrent falls
weighted 1. One additional point is added per decade of age above 50.
Categories absent from the original Charlson instrument reflect the
condition groupings that could be resolved from the unmapped source
vocabulary during curation.

\section{Operational Definitions of Complications}
\label{app:definitions}

Definitions are given in the order used in the Results tables. As described
in Section~\ref{sec:methods}, the glycemic endpoints and acute kidney injury
are derived within the knowledge base from the corresponding measurements,
while ketoacidosis, infection, hyperosmolality and myocardial infarction were
computed during data preparation and supplied to the engine as primitive
concepts.

\paragraph{Thirty-day mortality.}
Death recorded in the linked mortality data between 14 days before and 30
days after the recorded discharge datetime. The backward tolerance ensures
that in-hospital deaths are captured where the death record precedes the
administrative discharge timestamp.

\paragraph{Hyperglycemia.}
Either a single glucose measurement $\geq250$ mg/dL, or a sustained
hyperglycemic state, defined as a glucose measurement $\geq180$ mg/dL whose
immediately preceding measurement within the previous 24 hours was also
$\geq180$ mg/dL. The sustained-state requirement prevents an isolated
elevated reading from being counted as a complication.

\paragraph{Severe hyperglycemia.}
Any single glucose measurement $\geq250$ mg/dL. This endpoint overlaps by
construction with hyperglycemia and is retained as a separate, more specific
target.

\paragraph{Hypoglycemia.}
Either a single glucose measurement $\leq54$ mg/dL, or a sustained
hypoglycemic state, defined as a glucose measurement $\leq70$ mg/dL whose
immediately preceding measurement within the previous 24 hours was also
$\leq70$ mg/dL.

\paragraph{Severe hypoglycemia.}
Any single glucose measurement $\leq54$ mg/dL.

\paragraph{Inpatient acute kidney injury (AKI).}
Laboratory-defined AKI of stage 2 or above, established either by a serum
creatinine ratio $\geq2.0$ relative to the patient's baseline, or by an
absolute serum creatinine $\geq4.0$ mg/dL. Stage~1 injury was deliberately
excluded from the endpoint.

\paragraph{Inpatient myocardial infarction (MI).}
Evidence of an ischemic cardiac event---ischemic heart disease, myocardial
infarction, coronary intervention, cardiac catheterization, or
revascularization---together with a troponin measurement of at least 600.
The troponin requirement excludes coded cardiac history unaccompanied by
biochemical evidence of acute injury.

\paragraph{Hyperosmolality.}
Emitted by either of two paths. The derived path implements a
hyperosmolar-hyperglycemic-state rule: for each glucose measurement
$\geq600$ mg/dL, the nearest serum sodium measurement within $\pm6$ hours is
located and the effective osmolality computed as
$2 \times \mathrm{Na} + \mathrm{glucose}/18$; an event is emitted where this
exceeds 320 mOsm/kg, provided no acidosis event occurs within $\pm6$ hours
of the triggering glucose draw. The acidosis exclusion separates
hyperosmolar states from diabetic ketoacidosis. The observation path
additionally emits an event wherever hyperosmolality is directly coded.

\paragraph{Ketoacidosis.}
Emitted by either of two paths. The observation path emits an event wherever
ketoacidosis is directly coded. The derived path requires the conjunction of
a pH measurement $\leq7.3$, a bicarbonate measurement $\leq10$ mEq/L within
24 hours of it, and intravenous insulin administration within 6 hours of
it---that is, biochemical acidosis together with the treatment response it
implies.

\paragraph{Infection.}
Requires a physiological trigger, corroborating diagnostic evidence, and
treatment, all within a single episode. The trigger is any one of: fever
onset, defined as at least two body-temperature measurements $\geq37.8^\circ$C
where all temperatures in the preceding 48 hours were $\leq37.7^\circ$C;
leukocytosis, defined as a white-cell count $\geq11$ where counts in the
preceding 48 hours lay between 4 and 11; leukopenia, defined as a white-cell
count $\leq4$ under the same prior-stability condition; or neutropenia,
defined as a neutrophil count $\leq1.5$ where preceding counts within 48
hours exceeded 1.5. Diagnostic evidence requires at least one blood culture,
urine culture, or chest radiograph within $\pm24$ hours of the trigger, the
latest such record within the window serving as the episode anchor.
Treatment requires at least one antibiotic administration after that anchor
and within 72 hours of the trigger. Triggers occurring within 72 hours of an
already-accepted event are collapsed into it, so an admission may contribute
multiple distinct episodes. The prior-stability condition attached to every
trigger restricts the endpoint to infections arising during the
hospitalization rather than present on admission.

\section{Assessed Guideline Recommendations}
\label{app:patterns}

Fourteen quality-assessment patterns were evaluated, corresponding to the
components reported in the Results
(Table~\ref{tab:pattern-definitions}). Each pattern defines an applicability
condition, so that it is scored only for admissions and time points at which
the underlying recommendation applies, and one or more compliance aspects.
Three aspect types are used: \emph{time}, whether a required action occurred
within a required interval; \emph{cyclic}, whether a recurring action
occurred at the required frequency; and \emph{value}, whether an
administered dose met the required magnitude. Insulin administration after
hyperglycemia is a limiting case of the time aspect: it is defined by an
overlap relation rather than by a delay, so no distance separates anchor
from action for a membership function to grade, and the constraint is
satisfied on existence alone---scoring 1 where the action is present and 0
where it is not.

\begin{table}[t]
\centering
\caption{Assessed guideline recommendations and their operational
definitions in the knowledge base.}
\label{tab:pattern-definitions}
\small
\begin{tabularx}{\textwidth}{
>{\raggedright\arraybackslash}p{0.30\textwidth}
l
>{\raggedright\arraybackslash}X
}
\toprule
Recommendation & Aspect & Operational definition \\
\midrule

Admission glucose monitoring & Cyclic &
Glucose measured in each 24-hour window over the first days of admission:
three days where diabetes is documented, two otherwise. Full credit at
three or more measurements per window, decaying linearly to zero at
none. \\

Routine glucose monitoring & Cyclic &
Glucose measured in each 24-hour window throughout the admission (evaluated
to 14 days). While glucose is in a hyperglycemic or hypoglycemic state, full
credit is given at three or more measurements per window. Under routine
conditions the trapezoid peaks at at least one measurement per window. \\

Glucose monitoring during corticosteroid therapy & Cyclic &
Glucose measured in each 24-hour window over the first two days of
corticosteroid treatment, with full credit at at least three measurements
per window. \\

BMI measurement on admission & Time &
Body mass index or weight recorded within 48 hours of admission, decaying
linearly to zero at 72 hours; a measurement recorded up to six months
before admission is carried forward as relevant. \\

Creatinine measurement on admission & Time &
Serum creatinine measured within 48 hours of admission, decaying linearly
to zero at 72 hours. Unlike BMI, no pre-admission measurement is carried
forward, so the measurement must fall within the admission itself. \\

Antidiabetic discontinuation on admission & Time &
Prescribed antidiabetic agents carrying high hypoglycemia risk discontinued
during the first 72 hours of admission, with credit decreasing linearly to
zero if an agent was administered during that time. \\

SGLT2 discontinuation on admission & Time &
Prescribed SGLT2 inhibitors discontinued during admission. Any in-hospital
administration scores 0, and the absence of one scores 1, assuming the
patient was treated with the agent at home. \\

Basal insulin continuation on admission & Time &
Basal insulin administered within 36 hours of admission for patients
established on basal insulin at home, decaying linearly to zero at 48
hours. \\

Insulin administration after hyperglycemia & Time &
Any dose of reactive bolus or intravenous insulin administered overlapping a
meal (2-hour window before or after the meal) while glucose is elevated. \\

Antidiabetic discontinuation after hypoglycemia & Time &
Prescribed antidiabetic agents withheld following a hypoglycemic event,
with credit increasing over 72 hours and full credit sustained
thereafter. \\

Basal dose reduction after hypoglycemia & Value &
Basal insulin dose reduced relative to the preceding dose following a
hypoglycemic event, scored on the ratio of the new dose to the preceding
one: any reduction receives full credit, declining to zero at an unchanged
dose. Unlike the bolus pattern below, hypoglycemic events \emph{not}
followed by a basal dose within 24 hours are scored as compliant rather than
excluded, following the convention that a withheld action satisfies the
recommendation. \\

Bolus dose reduction after hypoglycemia & Value &
Bolus insulin dose reduced following a hypoglycemic event, on the same dose
ratio and membership function as the basal pattern above, but evaluated only
where bolus insulin was actually administered within 24 hours of the event;
events with no subsequent bolus are excluded from scoring rather than
counted as compliant (assuming basal is routine and bolus is a reaction to
an event). \\

Routine basal dosage (frequency) & Cyclic &
Basal insulin administered in each 24-hour window through the first 14 days
of admission, evaluated in windows where hyperglycemia is present, with full
credit between one and four administrations per window and declining credit
beyond that range. \\

Routine basal dosage (dose) & Value &
Judged against body weight for patients not established on basal insulin at
home: full credit between 0.2 and 0.6 units per kilogram, rising from zero
below 0.2 and declining to zero at 1.0 units per kilogram, using a default
weight of 81~kg where no weight measurement is available. For patients
already receiving basal insulin at home, the established dose is accepted as
given and is not weight-normalized, assuming medical staff reacted to this
dose and might have increased the hospital dose due to tolerance. \\

\bottomrule
\end{tabularx}
\end{table}
Each aspect is scored by a trapezoidal membership function as described in
Section~\ref{sec:methods}, so the reported values are continuous on $[0,1]$:
full credit is given throughout the interval in which the recommendation is
considered met, with linear partial credit on either side of it and zero
beyond the outer bounds. Several trapezoids are deliberately open on one
side---full credit at or above a required count, for instance, with no
penalty for exceeding it---so that partial credit applies only where
departing from the recommendation in that direction is meaningful. A
recommendation carrying more than one aspect is reported under each aspect
separately, which is why routine basal dosage appears twice above---once for
the frequency of administration and once for the dose administered.

The differing treatment of unanswered anchors in the two dose-reduction
patterns is deliberate and reflects how the two insulins are used. Basal
insulin is given routinely, so a hypoglycemic event with no basal dose in
the following 24 hours represents the dose having been withheld, which
satisfies the recommendation and scores 1. Bolus insulin is given
reactively, so the absence of a dose carries no information about whether a
reduction was warranted, and those events are excluded from scoring rather
than counted as compliant. This is the principal reason the two components
report such different means ($0.87$ and $0.28$) despite sharing a membership
function.

The knowledge base additionally contains a metformin continuation pattern.
It is defined for detection only and carries no compliance function, so it
yields no score and is not among the fourteen assessed components.
